
\documentclass{article}
\usepackage{colortbl}
\usepackage{makecell}
\usepackage{multirow}
\usepackage{supertabular}

\begin{document}

\newcounter{utterance}

\centering \large Interaction Transcript for game `wordle\_withclue', experiment `medium\_frequency\_words\_clue\_no\_critic', episode 18 with reflexion{-}Qwen3.5{-}9B.
\vspace{24pt}

{ \footnotesize  \setcounter{utterance}{1}
\setlength{\tabcolsep}{0pt}
\begin{supertabular}{c@{$\;$}|p{.15\linewidth}@{}p{.15\linewidth}p{.15\linewidth}p{.15\linewidth}p{.15\linewidth}p{.15\linewidth}}
    \# & \multicolumn{2}{c}{Player} && \multicolumn{2}{c}{Game Master} \\
    \hline

    \theutterance \stepcounter{utterance}  
    & & & \multicolumn{4}{p{0.6\linewidth}}{
        \cellcolor[rgb]{0.9,0.9,0.9}{
            \makecell[{{p{\linewidth}}}]{
                \texttt{\tiny{[P1$\langle$GM]}}
                \texttt{You have attempted this same task before and did not succeed. These are the plans you wrote after those attempts.} \\
\texttt{Trial \#1: New plan: The last attempt was aborted because the response contained excessive text, a long list of brainstormed words, and incomplete sentences, violating the strict format requirement to provide only a concise explanation followed immediately by the guess. Additionally, the explanation failed to explicitly integrate the clue "court activity" with the feedback logic. Next time, I will strictly adhere to the format: a single, concise sentence in the "explanation:" tag that references the clue and the specific feedback (e.g., "Since 'i' is the only correct letter in the third position for 'court activity', I will try a word like 'judge' or 'order'"), followed immediately by a valid five{-}letter guess in the "guess:" tag without any extra commentary or lists.} \\
\\ 
\texttt{Trial \#2: New plan: The last attempt was aborted because the response contained excessive brainstorming, repetitive self{-}correction, and incomplete sentences, violating the strict format requirement to provide only a concise explanation followed immediately by a valid five{-}letter guess. Additionally, the explanation failed to explicitly integrate the clue "court activity" with the specific feedback logic (u=2nd, g=not 4th). Next time, I will strictly adhere to the format: a single, concise sentence in the "explanation:" tag that references the clue "court activity" and the specific feedback (e.g., "Since 'u' is confirmed in the second position and 'g' is elsewhere for 'court activity', I will try 'gavel' again or a variant like 'gavel' without 'v','e','l'"), followed immediately by a valid five{-}letter guess in the "guess:" tag without any extra commentary, lists, or rambling thoughts.} \\
\\ 
\\ 
\texttt{You are a language wizard who likes to guess words by using the given rules.} \\
\\ 
\texttt{Welcome to Wordle! You have six attempts to guess the target word, a valid English word of five lowercase letters (a{-}z). Please use the tags "explanation:" and "guess:" to provide a concise explanation for each guess.} \\
\\ 
\texttt{To help you make an informed guess, you will receive a clue for the word, such as} \\
\texttt{clue: snowy white.} \\
\\ 
\texttt{Here is an example guess based on the clue:} \\
\texttt{explanation: In the fairy tail Snow White, the girl is killed because she eats a poisoned apple. And the word apple has 5 letters.} \\
\texttt{guess: apple} \\
\\ 
\texttt{After each guess, your answer will be validated, and you will receive feedback indicating which letters are correct (green), which letters are correct but in the wrong position (yellow), and which letters are incorrect (red). This feedback can be useful in determining which letters to include or exclude in your next guess.} \\
\\ 
\texttt{For example, the feedback for "apple" might be:} \\
\texttt{guess\_feedback: a<yellow> p<yellow> p<green> l<yellow> e<red>} \\
\\ 
\texttt{The explanation should contain details about how both the guess\_feedback and clue work together in tandem to formulate a new guess.} \\
\\ 
\texttt{Do not generate any other text.} \\
\\ 
\texttt{Let's begin with your first guess.} \\
\\ 
\texttt{clue: court activity} \\
            }
        }
    }
     \\ \\

    \theutterance \stepcounter{utterance}  
    & \multicolumn{4}{p{0.6\linewidth}}{
        \cellcolor[rgb]{0.561,0.737,0.561}{
            \makecell[{{p{\linewidth}}}]{
                \texttt{\tiny{[P1$\rangle$GM]}}
                \texttt{explanation: Since the clue "court activity" suggests legal proceedings and the word must be five letters, I will start with "gavel" as a common term for a court tool.} \\
\texttt{guess: gavel} \\
            }
        }
    }
    & & \\ \\

    \theutterance \stepcounter{utterance}  
    & & & \multicolumn{2}{p{0.3\linewidth}}{
        \cellcolor[rgb]{0.9,0.9,0.9}{
            \makecell[{{p{\linewidth}}}]{
                \texttt{\tiny{[GM$|$GM]}}
                \texttt{attempts = 1} \\
\texttt{target\_word = suing} \\
\texttt{target\_word\_clue = court activity} \\
\texttt{guess = gavel} \\
\texttt{guess\_feedback = g<yellow> a<red> v<red> e<red> l<red>} \\
            }
        }
    }
    & & \\ \\

    \theutterance \stepcounter{utterance}  
    & & & \multicolumn{4}{p{0.6\linewidth}}{
        \cellcolor[rgb]{0.9,0.9,0.9}{
            \makecell[{{p{\linewidth}}}]{
                \texttt{\tiny{[P1$\langle$GM]}}
                \texttt{guess\_feedback: g<yellow> a<red> v<red> e<red> l<red>} \\
\\ 
\texttt{Provide your response only in this format.} \\
\texttt{explanation: details} \\
\texttt{guess: word} \\
            }
        }
    }
     \\ \\

    \theutterance \stepcounter{utterance}  
    & \multicolumn{4}{p{0.6\linewidth}}{
        \cellcolor[rgb]{0.561,0.737,0.561}{
            \makecell[{{p{\linewidth}}}]{
                \texttt{\tiny{[P1$\rangle$GM]}}
                \texttt{explanation: The feedback shows 'g' is correct but in the wrong position (not first), while 'a', 'v', 'e', 'l' are incorrect; since "court activity" often involves a judge presiding, I will try "judge" with 'g' moved to the end or a new word starting with 'j'.} \\
\texttt{guess: judge} \\
            }
        }
    }
    & & \\ \\

    \theutterance \stepcounter{utterance}  
    & & & \multicolumn{2}{p{0.3\linewidth}}{
        \cellcolor[rgb]{0.9,0.9,0.9}{
            \makecell[{{p{\linewidth}}}]{
                \texttt{\tiny{[GM$|$GM]}}
                \texttt{attempts = 2} \\
\texttt{target\_word = suing} \\
\texttt{target\_word\_clue = court activity} \\
\texttt{guess = judge} \\
\texttt{guess\_feedback = j<red> u<green> d<red> g<yellow> e<red>} \\
            }
        }
    }
    & & \\ \\

    \theutterance \stepcounter{utterance}  
    & & & \multicolumn{4}{p{0.6\linewidth}}{
        \cellcolor[rgb]{0.9,0.9,0.9}{
            \makecell[{{p{\linewidth}}}]{
                \texttt{\tiny{[P1$\langle$GM]}}
                \texttt{guess\_feedback: j<red> u<green> d<red> g<yellow> e<red>} \\
\\ 
\texttt{Provide your response only in this format.} \\
\texttt{explanation: details} \\
\texttt{guess: word} \\
            }
        }
    }
     \\ \\

    \theutterance \stepcounter{utterance}  
    & \multicolumn{4}{p{0.6\linewidth}}{
        \cellcolor[rgb]{0.561,0.737,0.561}{
            \makecell[{{p{\linewidth}}}]{
                \texttt{\tiny{[P1$\rangle$GM]}}
                \texttt{explanation: The feedback confirms 'u' is correct in the second position and 'g' is correct but not in the fourth position, while 'j', 'd', 'e' are incorrect; since "court activity" involves a judge speaking, I will try "gavel" again but place 'g' in the first position as it was yellow in the first guess, or try "gavel" with 'g' at the start. Wait, 'g' was yellow in the first guess meaning it's in the word but not first. In the second guess, 'g' was yellow in the fourth position, meaning it's in the word but not fourth. So 'g' is in the word, not first, not fourth. 'u' is second. Let's try "gavel" again? No, 'a','v','e','l' were red in first guess. So 'g' is the only known letter besides 'u'. Maybe "guru"? No. "guru" doesn't fit court. "gavel" has 'g','a','v','e','l'. 'a','v','e','l' red. So only 'g' and 'u' are in the word. 'g' not first, not fourth. 'u' second. So pattern: ? u ? g ?. Clue "court activity". Maybe "gavel" is wrong. What about "gavel" without those letters? Maybe "gavel" is not the word. What is a 5 letter court activity with u second and g not first not fourth? "gavel" is out. "judge" is out. "gavel" has g,a,v,e,l. All except g red. So only g and u are in. So word has u second, g somewhere not 1 or 4. Maybe "gavel" is not it. What about "gavel"? No. Maybe "gavel" is wrong. Let's think of court activities: trial, jury, verdict, gavel, judge, lawyer, witness, appeal, arrest, bail, bond, case, court, crime, debt, fine, guilt, jury, law, legal, order, plea, rule, suit, trial, verdict. 5 letters: trial, jury, verdict (7), gavel (5), judge (5), lawyer (6), witness (7), appeal (6), arrest (6), bail (4), bond (4), case (4), court (5), crime (5), debt (4), fine (4), guilt (5), jury (4), law (3), legal (5), order (5), plea (4), rule (4), suit (4), trial (5). So possible: trial, court, crime, guilt, legal, order, gavel, judge. From feedback: u is second. So ? u ? ? ?. Court} \\
            }
        }
    }
    & & \\ \\

    \theutterance \stepcounter{utterance}  
    & & & \multicolumn{2}{p{0.3\linewidth}}{
        \cellcolor[rgb]{0.9,0.9,0.9}{
            \makecell[{{p{\linewidth}}}]{
                \texttt{\tiny{[GM$|$GM]}}
                \texttt{Error: The response should contain the 'guess:' keyword exactly once.} \\
            }
        }
    }
    & & \\ \\

    \theutterance \stepcounter{utterance}  
    & & & \multicolumn{2}{p{0.3\linewidth}}{
        \cellcolor[rgb]{0.9,0.9,0.9}{
            \makecell[{{p{\linewidth}}}]{
                \texttt{\tiny{[GM$|$GM]}}
                \texttt{game\_result = ABORT} \\
            }
        }
    }
    & & \\ \\

    \theutterance \stepcounter{utterance}  
    & & \multicolumn{4}{p{0.6\linewidth}}{
        \cellcolor[rgb]{0.95,0.95,0.95}{
            \makecell[{{p{\linewidth}}}]{
                \texttt{\tiny{[P1$\leftrightarrow$P]}}
                \texttt{Post{-}episode Reflexion (not sent to GM):} \\
\texttt{New plan: The previous attempts were aborted because the response contained excessive internal monologue, brainstorming lists, and incomplete sentences, violating the strict format requirement. Next time, I will strictly adhere to the format by providing only a single, concise sentence in the "explanation:" tag that explicitly integrates the clue "court activity" with the specific feedback (u in position 2, g not in position 1 or 4), followed immediately by a valid five{-}letter guess in the "guess:" tag without any extra commentary, lists, or rambling thoughts.} \\
            }
        }
    }
    & \\ \\

\end{supertabular}
}

\end{document}
